\chapter{Equilibrium, local models and the six controllers}
\section{Keep the full model as the prediction model}
Collect the nonlinear mechanical, rotor and electrical states:
\begin{equation}
\vect x=(\vect r^T,\vect v^T,q^T,\vect\omega^T,
\Omega_1,\ldots,\Omega_4,I_1,\ldots,I_4,V_p,z)^T.
\end{equation}
The preceding equations define $\dot{\vect x}=f(\vect x,\vect u,\vect\theta)$,
where $\vect\theta$ contains measured or assumed physical parameters. The unit
quaternion has one constraint. A numerical simulation integrates these equations;
it need not replace trigonometric functions with their small-angle forms.

For controller calculation, select a level hover operating point:
\begin{align}
\vect v_0&=0,&\vect\omega_0&=0,&R_0&=I,\\
\sum_i k_T\Omega_{i0}^2&=mg,&
k_\tau I_{i0}&=Q_i(\Omega_{i0})+B_r\Omega_{i0}.
\end{align}
For symmetric rotors, $\Omega_{i0}=\sqrt{mg/(4k_T)}$.
Without identified $k_T$, this is a relationship, not an available numeric RPM.

\section{Where a local linear model comes from}
Write $\vect x=\vect x_0+\delta\vect x$ and $\vect u=\vect u_0+\delta\vect u$.
A first-order Taylor expansion at equilibrium gives
\begin{align}
f(\vect x_0+\delta\vect x,\vect u_0+\delta\vect u)
&=f(\vect x_0,\vect u_0)
+A\delta\vect x+B\delta\vect u+\text{higher-order terms},\\
A&=\left.\frac{\partial f}{\partial\vect x}\right|_0,&
B&=\left.\frac{\partial f}{\partial\vect u}\right|_0.
\end{align}
At equilibrium $f(\vect x_0,\vect u_0)=0$. Neglecting only the remainder gives
the \emph{local} perturbation model. This step is explicitly labelled and does
not delete the original nonlinear equations.

For example, speed-squared thrust expands exactly as
\begin{align}
T(\Omega_0+\delta\Omega)
&=k_T\Omega_0^2+2k_T\Omega_0\delta\Omega+k_T(\delta\Omega)^2,\\
\delta T&\approx 2k_T\Omega_0\delta\Omega.
\end{align}
The local slope $2k_T\Omega_0$ changes with operating point.
For a measured command curve $T=f_T(u)$,
$\delta T\approx f_T'(u_0)\delta u$. A nonlinear thrust curve can therefore be
used for simulation while its measured hover slope supplies local controller gain.

\section{The currently reproducible baseline}
\begin{center}
\begin{tabular}{@{}ll@{}}\toprule
Parameter & Assumed baseline\\\midrule
$m$ & 1.2 kg\\
$(J_x,J_y,J_z)$ & (.012,.012,.022) kg m$^2$\\
$D$, $a=D/(2\sqrt2)$ & .45 m, .159099 m\\
$T_{\max}$ & 6 N per rotor in the normalized linear-thrust surrogate\\
$\tau_m$ & .04 s thrust lag\\
$d$ & .12 N s/m\\
$\Delta t$ & .002 s virtual control interval\\\bottomrule
\end{tabular}
\end{center}
For the surrogate $T_i=T_{\max}u_i^a$,
\begin{equation}
c_0=\frac{mg}{4T_{\max}}=.4903325.
\end{equation}
The other model is the provisional .740 kg component URDF, with inertia
\begin{equation}
(J_x,J_y,J_z)=(.00885852,.00906322,.01718682)\unit{kg\,m^2}.
\end{equation}
Its masses are not verified and it is
not the physics body of the recorded 1.2 kg flight. Never combine one model's
mass with another model's inertia accidentally.

\section{Derive each local transfer function}
For a small roll perturbation, with other rates zero,
\begin{align}
J_x\dot p&=4aT_{\max}u_\phi^a,\\
\tau_m\dot u_\phi^a+u_\phi^a&=u_\phi.
\end{align}
In radian units define $b_x=4aT_{\max}/J_x$. Under zero perturbation initial
conditions, $\mathcal L\{\dot p\}=sP(s)$, so
\begin{align}
sP(s)&=b_xU_\phi^a(s),\\
(\tau_ms+1)U_\phi^a(s)&=U_\phi(s),\\
\frac{P(s)}{U_\phi(s)}&=\frac{b_x}{s(\tau_ms+1)}.
\end{align}
Integrating rate gives a further factor $1/s$ in the local angle plant:
\begin{equation}
G_\phi(s)=\frac{b_x}{s^2(\tau_ms+1)}.
\end{equation}
Pitch substitutes $J_y$. Yaw rate substitutes $b_z=4\kappa T_{\max}/J_z$ and
has no outer yaw-angle controller.

The present C rate inputs are degrees/s. Multiply the radian effectiveness by
$180/\pi$ before using those C gain units:
\begin{equation}
b_{\phi,\deg}=18231.4054,\qquad b_{\psi,\deg}=937.5673.
\end{equation}
For height, subtract hover equilibrium and set $\lambda=d/m$:
\begin{align}
\delta\ddot h+\lambda\delta\dot h&=b_h\delta c^a,&b_h&=\frac{4T_{\max}}m,\\
G_h(s)&=\frac{b_h}{s(s+\lambda)(\tau_ms+1)}.
\end{align}
The numeric baseline is
$G_h(s)=20/[s(s+.1)(.04s+1)]$.

\section{Why do not all controllers use nonzero P, I and D?}
The letters name three available mathematical terms:
\begin{align}
e(t)&=r(t)-y(t),\\
u_P(t)&=K_pe(t),\\
u_I(t)&=K_i\int_0^t e(\xi)\dd\xi,\\
u_D(t)&=-K_d\dot y(t).
\end{align}
The last form differentiates measurement to avoid target-derivative kick.
Choosing $K_i=0$ or $K_d=0$ is a design choice, not an incomplete calculation.

\subsection{P changes the immediate response}
Larger error gives larger correction. For an angle error of 3 degrees and outer
$K_a=3\unit{s^{-1}}$, the desired rate is 9 degrees/s. P-only angle control can
work with an inner rate PI that rejects constant torque disturbances.

\subsection{I adds a state and removes certain steady errors}
For a constant rate disturbance, the inner rate integral can grow until its
output balances the disturbance torque. But each integrator adds dynamic order
and can accumulate unavailable demand during saturation. An additional outer
angle integral is not automatically beneficial; sensor bias is not made physically
correct merely by integrating its mistaken measurement.

\subsection{D amplifies high-frequency measurement variation}
If independent measurement noise has variance $\sigma_y^2$,
\begin{align}
\dot y_k&\approx(y_k-y_{k-1})/\Delta t,\\
\Var(\dot n_k)&=\frac{2\sigma_y^2}{\Delta t^2}.
\end{align}
At .2 degrees/s gyro sigma and .002 s spacing, the unfiltered rate derivative
noise sigma would be about $141.4$ degrees/s$^2$. Real filtering/correlation
changes that number, but not the need to account for noise. A nonzero rate D
term therefore needs bandwidth/noise design, not merely a wish to use all letters.

\subsection{Why altitude does use D}
The altitude controller has a useful filtered velocity estimate. It forms
\begin{equation}
u_h=K_p(h_{\rm ref}-\hat h)+K_iI_h+K_d(v_{\rm cmd}-\hat v).
\end{equation}
Since $\dot h_{\rm ref}\approx v_{\rm cmd}$ away from its limiter, the last term
is a derivative of height error using an estimated state rather than a noisy
finite difference. There is one altitude PID, not a seventh separate velocity PID.

\section{Calculate the initial six gains}
For a rate plant first approximated by $b/s$, use PI:
\begin{align}
C(s)&=\frac{K_ps+K_i}{s},\\
1+C(s)G(s)&=0,\\
s^2+bK_ps+bK_i&=0.
\end{align}
Choose $s^2+2\zeta\omega_ns+\omega_n^2$ and equate coefficients:
\begin{equation}
K_p=\frac{2\zeta\omega_n}{b},\qquad
K_i=\frac{\omega_n^2}{b}.
\end{equation}
For the current candidate, $\zeta=.9$, roll/pitch $\omega_n=8.4$ rad/s,
yaw $\omega_n=2.8$ rad/s, and outer angle $K_a=3$.
These bandwidths are chosen response objectives, not motor measurements.

For the initial altitude double-integrator design,
\begin{align}
G(s)&=b_h/s^2,\\
1+\left(K_p+K_i/s+K_ds\right)G(s)&=0,\\
s^3+b_hK_ds^2+b_hK_ps+b_hK_i&=0.
\end{align}
Expand desired $(s+p_h)(s^2+2\zeta_h\omega_hs+\omega_h^2)$:
\begin{align}
s^3+(2\zeta_h\omega_h+p_h)s^2
+(\omega_h^2+2\zeta_h\omega_hp_h)s+p_h\omega_h^2.
\end{align}
Coefficient matching gives
\begin{equation}
K_d=\frac{2\zeta_h\omega_h+p_h}{b_h},\quad
K_p=\frac{\omega_h^2+2\zeta_h\omega_hp_h}{b_h},\quad
K_i=\frac{p_h\omega_h^2}{b_h}.
\end{equation}
Use $\omega_h=1.05$, $p_h=.56$, $\zeta_h=.9$. The resulting degree-based C gains:
\begin{center}
\begin{tabular}{@{}lrrr@{}}\toprule
Loop & $K_p$ & $K_i$ & $K_d$\\\midrule
Roll angle & 3 & 0 & 0\\
Roll rate & .0008293381 & .0038702447 & 0\\
Pitch angle & 3 & 0 & 0\\
Pitch rate & .0008293381 & .0038702447 & 0\\
Yaw rate & .0053756141 & .0083620664 & 0\\
Altitude & .108045 & .03087 & .1225\\\bottomrule
\end{tabular}
\end{center}

\section{Now restore lag, drag and the complete cascade}
For a rate PI with lag, the characteristic is
\begin{equation}
\tau_ms^3+s^2+bK_ps+bK_i.
\end{equation}
Positive gains are not enough: the cubic Routh condition also requires
$K_p>\tau_mK_i$. The baseline roll/pitch polynomial is
$.04s^3+s^2+15.12s+70.56$, with poles $-7$ and $-9\pm13.0767j$.

The outer angle loop surrounds the closed inner rate loop plus angle integration:
\begin{align}
T_p(s)&=\frac{b(K_ps+K_i)}{\tau_ms^3+s^2+bK_ps+bK_i},\\
1+\frac{K_aT_p(s)}s&=0.
\end{align}
Multiplying out gives the complete cascade characteristic
\begin{equation}
\tau_ms^4+s^3+bK_ps^2+b(K_i+K_aK_p)s+K_abK_i.
\end{equation}
For altitude with drag and lag:
\begin{equation}
\tau_ms^4+(1+\tau_m\lambda)s^3+(\lambda+b_hK_d)s^2+b_hK_ps+b_hK_i.
\end{equation}
These quartics retain more of the local dynamics than the initial simplified
coefficient-matching models.

For $a_4s^4+a_3s^3+a_2s^2+a_1s+a_0$, the Routh first-column construction is
\begin{align}
b_1&=\frac{a_3a_2-a_4a_1}{a_3},\\
c_1&=\frac{b_1a_1-a_3a_0}{b_1}.
\end{align}
With all $a_i>0$, positivity reduces to
\begin{align}
\Delta_2&=a_3a_2-a_4a_1>0,\\
\Delta_3&=a_3a_2a_1-a_4a_1^2-a_3^2a_0>0.
\end{align}
For the baseline angle cascade these are 10.4832 and 1003.532544.
For altitude they are 2.473764 and 4.7232075492. Both local continuous models
satisfy the conditions; this does not prove a delayed/saturated nonlinear flight.

\subsection{A fuller altitude coefficient match}
Choose the fourth-order desired polynomial
\begin{equation}
\tau_m(s+\rho)(s+p_h)(s^2+2\zeta_h\omega_hs+\omega_h^2).
\end{equation}
The cubic coefficient fixes
$\rho=1/\tau_m+\lambda-p_h-2\zeta_h\omega_h$. Equate the other coefficients:
\begin{align}
K_d&=\frac{\tau_m[\omega_h^2+2\zeta_h\omega_hp_h+
\rho(p_h+2\zeta_h\omega_h)]-\lambda}{b_h},\\
K_p&=\frac{\tau_m[p_h\omega_h^2+
\rho(\omega_h^2+2\zeta_h\omega_hp_h)]}{b_h},\\
K_i&=\frac{\tau_m\rho p_h\omega_h^2}{b_h}.
\end{align}
This gives the uninstalled analytic candidate
$(.09912357,.02796822,.1103068)$ for $(K_p,K_i,K_d)$.
It has selected local poles $-22.65$, $-.56$, $-.945\pm.457684j$.
It has not replaced the recorded-flight gains or been verified in a new flight.

\section{Sampling, limits and nonlinear assessment}
For a local continuous state model, zero-order-hold discretization is
\begin{equation}
A_d=e^{A\Delta t},\qquad B_d=\int_0^{\Delta t}e^{A\xi}B\dd\xi.
\end{equation}
Augment the discrete plant with integral, estimator and delay states. Local
sampled stability requires all closed-loop eigenvalues to satisfy $|z_i|<1$.
Continuous poles alone do not prove this implementation property.

Use the nonlinear model to examine different initial attitudes, wind, changing
voltage, sensor errors, motor limits and parameter ranges. This is validation of
the design model, not a reason to linearize away every difficult term.
Saturation, ground contact and the battery disarm event are hybrid/non-smooth
events and must be represented explicitly.

\chapter{Kalman estimation: uncertainty changes the gains}
\section{From noisy sensors to states}
The nonlinear sensor model for a body accelerometer is
\begin{equation}
\vect a_m=R^T(\ddot{\vect r}-\vect g_W)+\vect b_a+\vect n_a,\qquad
\vect g_W=(0,0,-g)^T.
\end{equation}
At rest and level, it reads approximately $(0,0,g)^T$; in ideal free fall
$\ddot{\vect r}=\vect g_W$ and it reads approximately zero. An accelerometer
measures specific force, not a direct world acceleration vector.
The gyro model is $\vect g_m=\vect\omega+\vect b_g+\vect n_g$.

Gravity-based tilt observations are
\begin{equation}
\phi_a=\operatorname{atan2}(a_y,a_z),\qquad
\theta_a=\operatorname{atan2}(-a_x,\sqrt{a_y^2+a_z^2}).
\end{equation}
During translation these observations contain nongravitational acceleration.
Using a stationary noise variance does not remove that modelling error.

\section{The current two-state angle/bias filter}
For one independent angular axis:
\begin{align}
\vect x&=(\alpha,b_g)^T,&
F&=\begin{bmatrix}1&-\Delta t\\0&1\end{bmatrix},&
H&=[1,0],\\
\alpha^-&=\alpha+\Delta t(g_m-b_g),&
b_g^-&=b_g.
\end{align}
Although sometimes called a one-dimensional angle filter, it estimates two
states per axis. Let $P=\begin{bmatrix}a&b\\c&d\end{bmatrix}$. Then
\begin{equation}
FPF^T=
\begin{bmatrix}
a-\Delta t(b+c)+\Delta t^2d&b-\Delta t d\\
c-\Delta t d&d
\end{bmatrix}.
\end{equation}
The current noise model adds
\begin{equation}
Q=\begin{bmatrix}
\sigma_g^2\Delta t^2+q_b\Delta t^3/3&-q_b\Delta t^2/2\\
-q_b\Delta t^2/2&q_b\Delta t
\end{bmatrix}.
\end{equation}
Here $\sigma_g^2$ is rate variance per sample and $q_b$ is a continuous
bias-random-walk intensity. They do not have interchangeable units.

\section{Derive the correction quantities}
\begin{align}
P^-&=FP^+F^T+Q,\\
\nu&=z-H\vect x^-,\\
S&=HP^-H^T+R,\\
K&=P^-H^TS^{-1},\\
\vect x^+&=\vect x^-+K\nu.
\end{align}
For $H=[1,0]$, $S=P_{00}^-+R$ and
$K=(P_{00}^-,P_{10}^-)^T/S$. A scalar observation can correct bias through
cross covariance even though it does not directly measure bias.
The covariance uses Joseph form
\begin{equation}
P^+=(I-KH)P^-(I-KH)^T+KRK^T.
\end{equation}
With the existing assumed noise values and .002 s updates, limiting gains are
approximately $K_\alpha=.000966191$, $K_b=-.000223499$.
They are conditional on the model; startup gains differ and they are not
universal sensor constants.

\section{Height and vertical velocity}
For $\vect x=(h,v_z)^T$, external upward acceleration and interval $\Delta t$:
\begin{align}
F&=\begin{bmatrix}1&\Delta t\\0&1\end{bmatrix},&
G&=\begin{bmatrix}\Delta t^2/2\\\Delta t\end{bmatrix},\\
\vect x^-&=F\vect x^++Ga_{\rm up},&
Q&=G\sigma_a^2G^T.
\end{align}
With $\sigma_a=1.5$ m/s$^2$, $\Delta t=.002$ s:
\begin{equation}
Q=\begin{bmatrix}9\,10^{-12}&9\,10^{-9}\\9\,10^{-9}&9\,10^{-6}\end{bmatrix}.
\end{equation}
Predict at 500 Hz, correct only on independent fresh observations, typically
50 Hz for the assumed barometer. With $R_{\rm baro}=.0025$ m$^2$, periodic
limiting gains at correction time are approximately $(.0834306,.181649)^T$.
Applying these as constant gains every 2 ms would be a different filter.

An observed height from a tilted downward range sensor also needs its beam
geometry, sensor lever arm and terrain assumptions. A one-ray range is not
automatically aircraft vertical height over arbitrary obstacles. The current
simulation's floor ray is declared analytical, not a native physical sensor.

\section{A fuller estimator is possible, but must be identified}
For a nonlinear model $\dot x=f(x,u)$ and observations $y=h(x)$, an EKF uses
Jacobians of $f,h$ along the evolving estimate. A quaternion error-state filter
can include position, velocity, attitude and IMU biases. It preserves more of
the nonlinear observation geometry than independent tilt filters.
This is a future estimator option, not an implementation claim. Unknown noise,
unobservable states and poor camera/pose information do not disappear merely
because a larger filter is used.
